\subsection{Physical selection as a natural relation}
Away from the electronic center, a species must not select a single cell by decree. The twelve noncentral multisections are fibers stable under cellular inversion and do not, in general, admit an equivariant point section. The correct realization begins with a relation

\[
\mathscr S\subset \mathcal P\times\mathcal R_{324},
\qquad
\mathscr S(X)=\{\gamma:(X,\gamma)\in\mathscr S\},
\]
whose value may be a point, an orbit, or a fiber. For the electron, \(\mathscr S(e)=\{\gamma_e\}\) has rank one because \((5,5)\) is the unique fixed point. For color, flavor, and generation, \(\mathscr S(X)\) preserves the corresponding orbit until a physical coupling compresses it.

The relation is admissible if it satisfies simultaneously:

\begin{enumerate}
\item depends only on representation, orientation, memory, boundary, charge,
flavor, color, and generation constructed before the contrast;

\item commutes with particle--antiparticle conjugation;
\item is equivariant with respect to the chromatic action and cellular inversion;
\item is compatible with refinement, nonadic return, and dodecaphase reading;
\item preserves every multiplicity that has not been eliminated by a coupling;
\item remains invariant under permuting or replacing the external mass values.
\end{enumerate}
The sixth point is the decisive control. If altering the experimental column changes \(\mathscr S(X)\), the construction is retrospective. If the relation remains and only the contrast residual changes, the causal direction is preserved.

\Needspace{7\baselineskip}
